\documentclass{article}
\usepackage[a4paper, total={6.5in, 10in}]{geometry}
\setlength{\parindent}{0pt}
\usepackage{tcolorbox}
\tcbset{colback=white!94!black, colframe=black!60!white, coltitle = white, boxrule=0.8pt, arc=2mm}
\newtcolorbox{boxs}[2][]{title= #2,#1}
\usepackage{datetime2}
\usepackage[british]{babel}
\usepackage{amsfonts}
\usepackage{amsmath}

\title{\textbf{Probability Lecture 6}}
\author{Robert O'Connell}
\date{\today}
\begin{document}
\maketitle
\subsection*{Variance}
A measure of spread of a PMF
\\

\begin{itemize}
    \item Random Variable \(X\), with mean \(\mu = E[X]\)
    \item Distance from the mean: \(X - \mu\)
    \item Average Distance from the mean? \\
\(E[X-\mu] = E[X] - \mu = \mu - \mu = 0\) \dots not useful
\end{itemize}

Definition of Variance : \(\text{var}(X) = E[(X-\mu)^2]\)
\\

var\((X) = E[g(x)]\) where \(g(x) = (x-\mu)^2 = \sum_x (x-\mu)^2 p_X(x)\)
\\

Standard deviation: \(\sqrt{\text{var}(X)}\)

\subsection*{Properties of the Variance}
\[
\text{var}(aX+b) = a^2\text{var}(X)
\]

Let \(Y = X+b\)
\[
\nu = E[Y] = \mu + b
\]
\[
\text{var}(Y) = E[(Y-\nu)^2] = E[(x+b-(\mu +b))^2] = E[(x-\mu)^2] =\text{var}(X)
\]

Let \(Y = aX\)
\[
\nu = E[Y] = a\mu
\]
\[
\text{var}(Y) = E[(aX-a\mu)^2] = E[a^2(x-\mu)^2] = a^2E[(x-\mu)^2] = a^2\text{var}(X)
\]

\subsubsection*{Another Formula}
\[
\text{var}(X) = E[X^2] - (E[X])^2
\]

We have
\[
\text{var}(X) = E[(x-\mu)^2] = E[X^2 - 2\mu X + \mu^2] = E[X^2] -2\mu E[X] + \mu^2
\]
but \(\mu = E[X]\)
\[
\Rightarrow E[X^2] - (E[X])^2
\]

\subsection*{Variance of the Bernoulli}
i.e. \(X = 1 , 0\) with probability \(p, 1-p\)

\[
\text{var}(X) = \sum_{x}(x-E[X])^2p_X(x) = (1-p)^2p + (0-p)^2(1-p) = p(1-p)
\]

or
\[
\text(var)(X) = E[X^2]- (E[x])^2 = p - p^2 = p(1-p)
\]
as \(X^2 = X\) for a Bernoulli PMF

\subsection*{Variance of the uniform}
\[
\text{var}(X) = E[X^2] - (E[X])^2 = \frac{1}{n-1} (0^2 + 1^2 + 2^2 + \dots + n^2) - (\frac{n}{2})^2 = \frac{1}{12}n(n+2)
\]
more generally sub in \(b-a = n\)

\textbf{Note}\\
Could add more formulae to calculate var(\(X\))

\subsection*{Conditional PMF and expectation, given an event}
Condition on an event \(A \Rightarrow\) use conditional probabilties

\[
p_X(x|A) = P(X=x|A) 
\]
\[
\sum_x p_{X|A}(x) = 1
\]
\[
E[X|A] = \sum_x xp_{X|A}(x)
\]
\[
E[g(X)|A] = \sum_x g(x)p_{X|A}(x)
\]
\[
\text{var}(X) = E[(X|A-E[X|A])^2]
\]

\subsection*{Total Expectation Theorem}
We know the total probability theorem
\[
P(B) = P(A_1)P(B|A_1)+\dots + P(A_n)P(B|A_n)
\]
where \(A_1 \dots A_n\) form a partition of the sample space

Then, for all \(x\)
\[
P_X(x) = P(A_1) p_{X|A_1}(x) + \dots + P(A_n)p_{X|A_n}(x)
\]
\[
\sum_x xP_X(x) = P(A_1) \sum_x x p_{X|A_1}(x) + \dots
\]
Then the total expectation theorem is
\[
E[X] = P(A_1)E[X|A_1] + \dots + P(A_n)E[X|A_n]
\]



\subsection*{Conditioning a geometric random variable}
\(X\): number of independent coin tosses until first head; \(P(H) = p\)

\[
p_X(k) = (1-p)^{k-1}p, \quad k = 1,2,\dots
\]

\textbf{Memorylessnes:} Number of remaining coin tosses conditioned on Tails in the first toss is geometric with parameter \(p\)

i.e. Conditioned on \(X>n\), \(X-n\) is geometric with parameter \(p\)


\[
P_{X-1|X>1}(3) = P(X-1=3| X>1) = P(T_2T_3H_4 | T_1) = P(T_2T_3H_4) = (1-p)2p = P_X(3)
\]

\[
P_{X-n|Xn}(k) = p_X(k)
\]

\subsection*{The mean of the geometric}
From the definition
\[
E[X] = \sum_{k=1}^{\infty} kp_X(k) = \sum_{k=1}^{\infty} k(1-p)^{k-1}p
\]

Or,
\[
E[X] = 1 + E[X-1] = 1 + p E[X-1|X=1] + (1-p)E[X-1|X>1] = 1 + 0 + (1-p)E[X]
\]
\[
E[X] = \frac{1}{p}
\]

\subsection*{Multiple Random Variables and joint PMFs}
\[
X:p_X, \quad Y:p_Y \quad \text{(marginal PMFs)}
\]
Joint PMF: \(P_{X,Y}(x,y) = P(X=x \ \text{and} \ Y=y)\)

\[
\sum_x \sum_y p_{X,Y}(x,y) = 1
\]
\[
p_X(x) = \sum_y p_{X,Y}(x,y)
\]
\[
p_Y(y) = \sum_x p_{X,Y}(x,y)
\]

\subsubsection*{More than two random variables}
\[
p_{X,Y,Z}(x,y,z) P(X=x, Y=y, Z=z)
\]
\[
\sum_x \sum_y \sum_z p_{X,Y,Z}(x,y,z) = 1
\]
\[
p_X(x) = \sum_y \sum_z p_{X,Y,Z}(x,y,z)
\]
\[
p_{X,Y}(x,y) = \sum_z P_{X,Y,Z}(x,y,z)
\]

\subsubsection*{Functions of multiple random variables}
\(Z = g(X,Y)\)

PMF: \(P_Z(z) = P(Z=z) = P(g(X,Y) = z) = \sum_{(x,y): g(x,y) =z } p_{X,Y}(x,y)\)

Expected Value Rule
\[
E[g(X,Y)] = \sum_x \sum_y g(x,y) p_{X,Y}(x,y)
\]


\subsubsection*{Linearity of Expectation}
\[
E[aX + b] = aE[x] + b
\]
For multiple random variables
\[
E[X+Y] = E[X] + E[Y]
\]

\[
E[X+Y] = E(g(X,Y)), \quad (g(x,y) = x+y)
\]
\[
= \sum_x \sum_y (x+y)p_{X,Y}(x,y) = \sum_x \sum_y xp_{X,Y}(x,y) + \sum_x \sum_y yp_{X,Y}(x,y)
\]
\[
\sum_x x \sum_y p_{X,Y}(x,y) + \dots = \sum_x p_X(x) + \sum_y y p_Y(y) = E[X] + E[Y]
\]

One can easily see that this formula holds for more random variables.

\subsection*{The mean of the binomial}
\(X\): binomial with parameters \(n,p\) 
 - number of successes in \(n\) independent trials

\[
E[X] = \sum_{k=0}^{n} k \binom{n}{k} p^k (1-p)^{n-k}
\]

To compute this, \\
\(X_i = 1\) if \(i\)th trial is a success, else \(X_i = 0\)

\[
X = X_1 + \dots + X_n
\]
\[
E[X] = E[X_1] + \dots + E[X_n] = np
\]

\end{document}